% 排版由未修改的 ThuThesis v7.7.1 管理：四边 3cm、20bp 正文行距、
% 标题层级、目录、页眉页码、图表公式、脚注与参考文献样式。
% 不在正文中用 \vspace 或 \fontsize 修补版式。
% math-font 会触发 unicode-math；先加载 mathtools 以保持兼容。
% 明确采用 mathtools 的上下方括号、unicode-math 的冒号符号。
% 这两条上游选择通知记入日志 Info；保留所有其他警告和错误。
\ExplSyntaxOn
\msg_redirect_name:nnn {unicode-math} {mathtools-overbracket} {info}
\msg_redirect_name:nnn {unicode-math} {mathtools-colon} {info}
\ExplSyntaxOff
\usepackage{mathtools}
\ifIMDTPrint
  \thusetup{output=print}
\else
  \thusetup{output=electronic}
\fi
\thusetup{
  math-font=xits,
  title={\ThesisTitle},
  author={\StudentName},
  student-id={\StudentID},
  department={\SchoolName},
  degree-category={电子信息硕士},
  professional-field={互动媒体设计与技术},
}
\usepackage{amsthm}
\usepackage{tabularx,longtable,threeparttable,multirow}
\usepackage{siunitx}
\usepackage{tikz}
\usetikzlibrary{arrows.meta,positioning}
\usepackage{algorithm}
\usepackage{algpseudocode}
\usepackage[sort&compress]{natbib}
\bibliographystyle{thuthesis-numeric}
\usepackage{xcolor}
\usepackage{hyperref}
\hypersetup{hidelinks,unicode,bookmarksnumbered=true,
  pdftitle={\ThesisTitle},pdfauthor={\StudentName}}
\usepackage[nameinlink,noabbrev]{cleveref}
\crefname{figure}{图}{图}
\crefname{table}{表}{表}
\crefname{section}{节}{节}
\crefname{algorithm}{算法}{算法}
\crefformat{chapter}{第#2#1#3章}
\crefformat{equation}{式（#2#1#3）}
\floatname{algorithm}{算法}
\renewcommand{\algorithmicrequire}{\textbf{输入：}}
\renewcommand{\algorithmicensure}{\textbf{输出：}}
\makeatletter
\@addtoreset{algorithm}{chapter}
\makeatother
\renewcommand{\thealgorithm}{\thechapter.\arabic{algorithm}}
\newcolumntype{Y}{>{\centering\arraybackslash}X}
\DeclarePairedDelimiter{\norm}{\lVert}{\rVert}
\newcommand{\dif}{\mathop{}\!\mathrm{d}}
\graphicspath{{figures/}}
% 提示与待办只在草稿模式显示；正式正文写在这些宏之外。
\newcommand{\WriteHere}[1]{\ifIMDTDraft{\color{black!55}#1}\par\fi}
\newcommand{\TODO}[1]{\ifIMDTDraft{\color{red}\textbf{待办：}#1}\par\fi}
% 防止模板短段落被拉伸成稀疏页面；不会改变正文固定行距。
\raggedbottom
